\section{Synthetic certification --- why the safeguards are necessary}

\subsection{Opening}
Before applying TOS-CMI to EEG, we construct synthetic controls in which the true leakage and task
information are known (Gaussian-mixture generators with a computable Bayes oracle). These controls are not
designed to maximize benchmark accuracy; they are designed to test whether each safety layer prevents a
specific failure mode --- first-moment blindness, geometry-only deletion, weak conditional task gates, and
underpowered certification. The main outcome of this line is not default deletion, but \textbf{calibrated refusal
when deletion cannot be certified}.

\subsection{Why geometry alone is not enough}
A selective-invariance method that only deletes ``domain-rich'' directions implicitly assumes (i) that those
directions can be \emph{found}, and (ii) that removing them is \emph{safe} for the task. Neither is automatic: a
direction can be domain-rich yet still carry conditional task information, and a deletion that is exactly
task-preserving \emph{in the geometry} ($R T = T$) can still raise conditional task risk. We therefore treat
localization and deletion as steps that must be \textbf{measured and certified}, and use the synthetic suite to
show that each unguarded step has a concrete failure mode.

\subsection{Score-Fisher detects leakage missed by first-moment scatter}
The candidate subspace is read from the conditional score-Fisher operator (\S2.2), which is second-order in
the log-likelihood scores. By construction it responds to covariance- and synergy-structured leakage, where
a first-moment (mean-scatter) statistic is identically blind: on a covariance-only generator the
mean-scatter direction is a no-op while the score-Fisher direction recovers the leakage subspace \claimtag{C1}. This
is why localization uses the score-Fisher operator rather than class-mean differences.

\subsection{Direct-sum deletion can still be conditionally unsafe}
On a synergy (``explaining-away'') generator, the domain-rich subspace is task-light by first- and
second-moment geometry, so a direct-sum projector deletes it without any \emph{apparent} task cost --- yet the
Bayes oracle shows the deletion removes real conditional task information ($\Delta_Y > 0$). Geometry alone would
accept this deletion. This is the concrete sense in which \textbf{direct-sum geometry is necessary but not
sufficient for conditional task safety} \claimtag{C2}, and it motivates an explicit conditional task-risk gate.

\subsection{Weak task gates unsafe-accept; log-ratio gates reduce the failure}
A naive task gate --- a small nested critic estimating the conditional task risk --- \emph{unsafe-accepts} such
deletions: its probe task-risk upper bound sits far below the true Bayes gap, so the deletion passes (Fig.~\ref{fig:synth}B). Replacing it with a cross-fitted one-step plug-in log-ratio estimator shrinks the gap between the
estimated bound and the oracle, and the gap continues to close with calibration sample size (Fig.~\ref{fig:synth}C); the
unsafe-accept occurs only in the small-sample, weak-critic regime \claimtag{C3}. The estimator is the bottleneck ---
not the geometry --- so the gate must use a calibrated estimator, not a convenience critic.

\subsection{Power certification turns unsafe deletion into refusal}
Even a calibrated estimator can be \emph{underpowered}: at a given sample size it may be unable to distinguish a
safe deletion from an unsafe one. We therefore attach a power certificate (a minimum-detectable-effect
floor) and require both task safety (an upper bound on $\Delta_Y$) and a domain-gain lower bound before deleting.
With this in place the certified gate produces \textbf{zero unsafe-accepts} across the synthetic phase diagram,
but is deliberately conservative --- it abstains or rejects in most cells (Fig.~\ref{fig:synth}D). Abstention is the intended
behavior when deletion cannot be certified, and is reported as such (not as a missing result).

\begin{figure}[t]
\centering
\includegraphics[width=\linewidth]{fig2_synthetic_cert.pdf}
\caption{\textbf{Synthetic certification: geometry alone is not safety.} Gaussian-mixture controls with a
known Bayes oracle. (A) Certification frontier: the gate fires on real leakage (oracle detection
$\approx$1 at the boundary) while the plug-in estimator is conservative, and neither fires under the
no-leakage null. (B) A weak nested critic \emph{unsafe-accepts} a conditionally-unsafe deletion --- its probe
task-risk upper bound sits far below the true Bayes gap (the accepted point, in red). (C) The estimator gap
(Bayes $-$ probe UCB) shrinks with calibration sample size; the unsafe-accept occurs only at the smallest
$n$. (D) With the cross-fitted plug-in estimator and a power floor, the certified gate produces \emph{zero}
unsafe-accepts but is conservative (it abstains or rejects), i.e.\ it defaults to refusal rather than weak
invariance.}
\label{fig:synth}
\end{figure}

\subsection{Takeaway}
\emph{The synthetic suite does not show that selective deletion should be enabled by default; it shows why a
selective-invariance method needs explicit measurement, conditional task-risk testing, and abstention.}

This is the discipline we carry into the EEG study (\S4): localize, test conditional task risk, and refuse
when a safe, useful deletion cannot be certified.
